% ================= 第 1 章: 引言与记号 =================

\section{引言与记号}

\subsection{\mbox{}}

设 $p$ 为素数, $K$ 为 $\Qp$ 的有限扩张.
继 $p$ 进群（如 $\GL_n(K)$）的 $p$ 进表示理论中的最新发展 (\cite{ref:44}, \cite{ref:43}, \cite{ref:45}, \cite{ref:22}, \cite{ref:25}, \cite{ref:26}, 等等) 之后, 出现了这样的问题: 将 $\GL_n(K)$ 的 $p$ 进表示"联系"到 $\Gal(\overline{\Qp}/K)$ 的 $n$ 维 $p$ 进表示 (例如 Fontaine 的潜在半稳定表示), 其精神在于一种可能的 Langlands 式局部对应.
由 \cite{ref:7} 和 \cite{ref:8} 中的例子开创 (但许多数学家早已暗示), 这一课题已经取得了一定数量的发展 (\cite{ref:9}, \cite{ref:18}, \cite{ref:3}, \cite{ref:19}), 并获得了"$p$ 进局部 Langlands 对应"这一通称.
不幸的是, 这个新对应比它的 $\ell$-进局部"大姊"困难得多, 因此目前获得的非平凡结果 (甚至猜想) 全都限于 $\GL_2$, 甚至 $\GL_2(\Qp)$.

在本文中, 我们谦逊地拟探索迄今悬而未决的情形: $\GL_2(\Qp)$ 的 $p$ 进表示与 $\Gal(\overline{\Qp}/\Qp)$ 的二维可约 (且可能分裂) 潜在晶体表示相对应.
我们为这些潜在晶体表示中的大多数定义了这样的 $\GL_2(\Qp)$ 表示.
然后我们证明, 当 Galois 表示来自一个模形式时, 相应的 $\GL_2(\Qp)$ 表示自然地 (连同 Galois 表示) 出现在模曲线的完备平展上同调中.
这正是本文的主要"局部-整体相容性"结果.
具体地说, 就是要证明: 可以在上同调中从 $\GL_2(\Qp)$ 一侧检测与所考虑的模形式相关联的 $\Gal(\Q/\Q)$ 表示在 $p$ 处是否分裂.
为此我们使用两个要素: 其一是 $p$ 进比较定理 (用于与模形式相关联的 Galois 表示), 其二是我们其中一人的 Jacquet 函子的 $p$ 进理论.

现在更精确地描述本文的内容.
设 $\sigma_p \simeq \begin{pmatrix} \eta_1 & * \\ 0 & \eta_2 \varepsilon^{-1} \end{pmatrix}$ 是 $\Gal(\overline{\Qp}/\Qp)$ 的可约潜在晶体表示, Hodge-Tate 权为 $(1-k,0)$, 其中 $k \ge 2$ 为整数 (即 $\eta_1$ 权为 $0$ 而 $\eta_2$ 权为 $2-k$), $\varepsilon$ 表示 $p$ 进分圆特征标.
再假设当 $k = 2$ 时 $\eta_1 \neq \eta_2$.
扩张 $*$ 若非零则唯一.
通过局部互反律把 $\Gal(\overline{\Qp}/\Qp)$ 的特征标看作 $\Qp^{\times}$ 的特征标, 我们如下给 $\sigma_p$ 配备一个带有 $\GL_2(\Qp)$ 连续酉作用的 $p$ 进 Banach 空间 $B(\sigma_p)$:

\medskip
\noindent (i) 若 $\sigma_p \simeq \begin{pmatrix} \eta_1 & 0 \\ 0 & \eta_2 \varepsilon^{-1} \end{pmatrix}$, 则
\[
B(\sigma_p) = \cind_{B(\Qp)}^{\GL_2(\Qp)} \eta_2 \otimes \eta_1 \;\oplus\; \cind_{B(\Qp)}^{\GL_2(\Qp)} \eta_1 \varepsilon \otimes \eta_2 \varepsilon^{-1},
\]
\noindent (ii) 若 $\sigma_p \simeq \begin{pmatrix} \eta_1 & * \\ 0 & \eta_2 \varepsilon^{-1} \end{pmatrix}$ 且 $* \neq 0$, 则 $B(\sigma_p)$ 是非分裂扩张:
\[
0 \to \cind_{B(\Qp)}^{\GL_2(\Qp)} \eta_2 \otimes \eta_1 \to B(\sigma_p) \to \cind_{B(\Qp)}^{\GL_2(\Qp)} \eta_1 \varepsilon \otimes \eta_2 \varepsilon^{-1} \to 0,
\]
其中记号 $\cind_{B(\Qp)}^{\GL_2(\Qp)} \eta \otimes \eta'$ 表示 $\GL_2(\Qp)$ 上取值于 $\Qp$ 的一个 (固定的) 有限扩张的连续函数 $f$, 满足 $f(bg) = (\eta \otimes \eta')(b)f(g)$ ($\GL_2(\Qp)$ 的作用是函数上通常的右平移).
情形 (ii) 中的 Banach 空间 $B(\sigma_p)$ 通过取局部解析抛物线诱导 $\anind_{B(\Qp)}^{\GL_2(\Qp)} \eta_1 |\,|^{k-1}\varepsilon^{2-k} \otimes \eta_2 |\,|^{1-k}\varepsilon^{k-2}$ 的唯一酉 $p$ 进完备化得到 (\cite{ref:44} 的意义), 其中 $|\,|$ 表示范数特征标 (参见 \ref{sec:2.2}).
因此上述对应的定义可以概括为这样一句话: "$\Gal(\overline{\Qp}/\Qp)$ 一侧分裂当且仅当 $\GL_2(\Qp)$ 一侧分裂".

\medskip
现在考虑模形式 $f = q + \sum_{n \ge 2} a_n(f)q^n$, 权 $k \ge 2$ 的归一化新尖点形式, 水平 $N$, 特征标 $\chi$, 是 Hecke 算子 $T_\ell$ ($(\ell,N) = 1$) 与 $U_\ell$ ($\ell \mid N$) 的特征向量.
记 $\sigma(f)$ 为与 $f$ 相关联的 $\Gal(\Q/\Q)$ 的 $p$ 进表示, $\sigma_p(f)$ 为它在 $p$ 处一个分解子群上的限制.
若 $M$ 是 $N$ 的与 $p$ 互素部分, $L$ 是 $f$ 定义在其上的 $\Qp$ 的有限扩张, 记 $\widehat{H}{}^{1}(K_1^{p}(M))_L := L \otimes_{\Zp}$ ($\Zp$-模 $\varinjlim_r H^{1}_{\mathrm{\'et}}(Y_1(M;p^r),\Zp)$ 的 $p$ 进完备化), 其中 $Y_1(M;p^r)$ 是与同余子群 $\Gamma_1(M) \cap \Gamma(p^r)$ 相关联的开模曲线.
我们定义 $\sigma(f)$-等型分量:
\[
\Pi_p(f) := \Hom_{\Gal(\Q/\Q)}\bigl(\sigma(f),\widehat{H}{}^{1}(K_1^{p}(M))_L\bigr),
\]
它自然地是一个带有 $\GL_2(\Qp)$ 连续酉作用的 $p$ 进 Banach 空间.
我们期望 $\GL_2(\Qp)$-表示 $\Pi_p(f)$ 精确地"包含"与 $p$ 进表示 $\sigma_p(f)$ 相同的"信息", 即包含形式 $f$ 的 $p$ 进 Hodge 理论 (参见例如 \cite{ref:9}).
现在假设 $\sigma_p(f)$ 是前述形式 (即潜在晶体且可约).
在这种情况下, 我们的局部-整体相容性猜想正是:

\begin{conjecture}\label{conj:1.1.1}
存在 $p$ 进 Banach 空间的拓扑 $\GL_2(\Qp)$-等变同构
\[
B(\sigma_p(f)) \simeq \Pi_p(f).
\]
\end{conjecture}

本文的主要结果是这个猜想的一个弱版本:

\begin{theorem}\label{thm:1.1.2}
\begin{enumerate}
\item[(i)] 恒有 $p$ 进 Banach 空间的闭 $\GL_2(\Qp)$-等变浸入:
\[
B(\sigma_p(f)) \hookrightarrow \Pi_p(f).
\]
\item[(ii)] 若 $\sigma_p(f)$ 不分裂, 则 (用前述记号) 还有:
\[
\Hom_{\GL_2(\Qp)}\bigl(B(\sigma_p(f)),\Pi_p(f)\bigr) = L,
\]
\[
\Hom_{\GL_2(\Qp)}\bigl(\cind_{B(\Qp)}^{\GL_2(\Qp)} \eta_1 \varepsilon \otimes \eta_2 \varepsilon^{-1},\Pi_p(f)\bigr) = 0.
\]
\end{enumerate}
\end{theorem}

我们现在说明证明定理 \ref{thm:1.1.2} 的步骤.
为简化, 我们在引言余下部分假设 $f$ 使得 $a_p(f)$ 是 $p$ 进单位 (对固定的嵌入 $\Q \hookrightarrow \Qp$).
若 $N$ 与 $p$ 互素, 记 $\alpha_p$ 和 $\beta_p$ 为 $X^2 - a_p(f)X + p^{k-1}\chi(p)$ 的根, 其中 $\beta_p$ 是 $p$ 进单位, 并设 $\tilde{f} = f(z) - \beta_p f(pz)$.
若 $p^r$ 恰整除 $N$ 且 $r > 0$, 设 $\tilde{f} = f\,|\,w_{p^r}$, 其中 $w_{p^r}$ 是 Atkin 算子 (参见 \ref{sec:4.1}).
回顾算子 $\theta$ 在 $q$-展开上表示 $q\,d/dq$.

\begin{theorem}\label{thm:1.1.3}
表示 $\sigma_p(f)$ 分裂当且仅当存在超收敛形式 $g$ (必为零斜率且权为 $2-k$) 使得
\[
\tilde{f} = \theta^{k-1}(g).
\]
\end{theorem}

注意方向 $\tilde{f} = \theta^{k-1}(g) \Rightarrow \sigma_p(f)$ 分裂较容易 (参见例如 \cite[Prop.~11]{ref:30}).
$g$ 的存在性更为微妙, 基于通常的 $p$ 进比较定理结合超收敛模形式理论.
形式 $g$ (若存在) 配得上 $f$ 的超收敛伴随形式之名, 因为定理 \ref{thm:1.1.3} 是 Gross 著名定理 (\cite{ref:33}) 在特征 $0$ 的类似物.
注意两种情形 ($\sigma_p(f)$ 分裂或不分裂) 在实践中确实都会出现 (例如, 若 $f$ 是 CM 的, 则 $g$ 由 \cite[Prop.~7.1]{ref:13} 恒存在, 故 $\sigma_p(f)$ 分裂).
定理 \ref{thm:1.1.3} 的某些情形已知 (通过把 \cite{ref:33} 的结果提升到特征 $0$ 的方法): 参见 \cite{ref:11} 以及 \cite[\S6]{ref:30}.

第二个结果的证明本质上使用了 \cite{ref:22} 和 \cite{ref:25} 中定义与研究的 Jacquet 函子的 $p$ 进版本.
称整权 $k \ge 2$ 的超收敛形式 $g$ 为"坏的", 若它不在算子 $\theta^{k-1}$ 的像中 ("坏"的精确定义见定义 \ref{def:5.4.1} 和命题 \ref{prop:5.4.4}).

\begin{theorem}\label{thm:1.1.4}
设 $g$ 是整权 $k \in \mathbb{Z}$、水平 $N$、特征标 $\chi$ 的超收敛模形式, 是 Hecke 算子的特征向量, 且 $U_p g = \alpha_p g$ 其中 $\alpha_p$ 非零.
再假设当 $k \ge 2$ 时 $g$ 不是"坏的".
则存在注入:
\[
\anind_{B(\Qp)}^{\GL_2(\Qp)} \nr(\alpha_p^{-1})\chi_p^{-1}\varepsilon^{2-k} \otimes \nr(\alpha_p) \hookrightarrow \widehat{H}{}^{1}(K_1^{p}(M))_L^{g},
\]
其中 $\widehat{H}{}^{1}(K_1^{p}(M))_L^{g}$ 是 $\widehat{H}{}^{1}(K_1^{p}(M))_L$ 关于 Hecke 算子 (在 $Np$ 之外) 作用按 $g$ 的特征值所取的特征空间, $\chi_p$ 是从 $\chi$ 得到的 $\Q$ 的有限 adele 特征标在 $p$ 处的分量, $\nr(x)$ 是 $\Qp^{\times}$ 的将 $p$ 送到 $x$ 的非分歧特征标, 而 "an" 表示局部解析抛物线诱导.
\end{theorem}

实际上, 文中证明了更精确的结果 (那里改用紧支上同调, 参见 \ref{sec:5.5}).
注意定理 \ref{thm:1.1.4} 中的注入因 $g$ 相关联的 (二维) Galois 表示出现在特征空间 $\widehat{H}{}^{1}(K_1^{p}(M))_L^{g}$ 中而不唯一.
定理 \ref{thm:1.1.4} 如下得到: 超收敛形式 $g$ 对应于 Coleman--Mazur Hecke 曲线 (或"本征曲线") (\cite{ref:17}) 上的一点.
根据 \cite{ref:24} 的理论, 若 $J_B\bigl(\widehat{H}{}^{1}(K_1^{p}(M))_L\bigr)$ 是对 $\widehat{H}{}^{1}(K_1^{p}(M))_L$ 应用 $p$ 进 Jacquet 函子 $J_B$ 后得到的 $\GL_2(\Qp)$ 的环面表示, 则与这点相关联的是 $J_B\bigl(\widehat{H}{}^{1}(K_1^{p}(M))_L\bigr)$ 关于环面与 Hecke 算子 (在 $Mp$ 之外) 作用的一个特征子空间: 在我们的情形, 它是特征标 $\nr(\alpha_p)|\,| \otimes \nr(\alpha_p^{-1})\chi_p^{-1}\varepsilon^{2-k}|\,|^{-1}$ 的特征子空间.
然后粗略地从函子 $J_B$ 的一个伴随律推导出定理 \ref{thm:1.1.4} (除一个本质情形需要更多工作外, 参见下文).

然后如下证明定理 \ref{thm:1.1.2}.
首先容易证明 $\cind_{B(\Qp)}^{\GL_2(\Qp)} \eta_2 \otimes \eta_1$ 这一块总包含在 $\Pi_p(f)$ 中.
假设 $\sigma_p(f)$ 分裂.
把定理 \ref{thm:1.1.4} 应用于定理 \ref{thm:1.1.3} 给出的超收敛形式 $g$, 并用 $\varepsilon^{k-1}$ 扭曲, 得到连续注入 $\anind_{B(\Qp)}^{\GL_2(\Qp)} \eta_1 \varepsilon \otimes \eta_2 \varepsilon^{-1} \hookrightarrow \Pi_p(f)$, 由此得到闭浸入 $\cind_{B(\Qp)}^{\GL_2(\Qp)} \eta_1 \varepsilon \otimes \eta_2 \varepsilon^{-1} \hookrightarrow \Pi_p(f)$.
因此在这种情况下 $\Pi_p(f)$ 包含两个连续抛物线诱导的直和, 即 $B(\sigma_p(f))$.
注意不能把定理 \ref{thm:1.1.4} 应用于超收敛形式 $\tilde{f}$, 因为此时 $\tilde{f}$ 恰好是"坏的".

现在假设 $\sigma_p(f)$ 不分裂.
则 $\Pi_p(f)$ 不能包含 $\cind_{B(\Qp)}^{\GL_2(\Qp)} \eta_1 \varepsilon \otimes \eta_2 \varepsilon^{-1}$ 这一块, 因为对前述推理取逆并使用 Hida 关于序超收敛形式的同调描述, 会得到 $\tilde{f}$ 在 $\theta^{k-1}$ 的像中, 从而蕴含 $\sigma_p(f)$ 分裂 (实际上, 我们证明关于 $\tilde{f}$ 的一个稍弱的断言就足以推出 $\sigma_p(f)$ 分裂, 参见定理 \ref{thm:5.7.2}).
但这次把定理 \ref{thm:1.1.4} 应用于形式 $\tilde{f}$ (它不再是"坏的") 给出连续注入:
\[
\anind_{B(\Qp)}^{\GL_2(\Qp)} \eta_1 |\,|^{k-1}\varepsilon^{2-k} \otimes \eta_2 |\,|^{1-k}\varepsilon^{k-2} \hookrightarrow \Pi_p(f),
\]
由此推出闭浸入 $B(\sigma_p(f)) \hookrightarrow \Pi_p(f)$, 因为按定义, $B(\sigma_p(f))$ 是局部解析表示 $\anind_{B(\Qp)}^{\GL_2(\Qp)} \eta_1 |\,|^{k-1}\varepsilon^{2-k} \otimes \eta_2 |\,|^{1-k}\varepsilon^{k-2}$ 的唯一酉 $p$ 进完备化.
再稍加工作, 得到定理 \ref{thm:1.1.2} 的 (ii) 中的重数 $1$.
注意在这种情况下定理 \ref{thm:1.1.4} 特别微妙, 因为我们处于"临界斜率"情形 (即超出 Amice--V\'elu 与 Vishik 定理的适用条件), 此情形的证明占据了第 5 部分的相当篇幅.

文本的划分如下: 在本引言和记号之后, 第 2 部分纯粹是局部的, 引入 $\GL_2(\Qp)$ 的各种抛物线诱导以及表示 $B(\sigma_p)$ 的构造.
第 3 部分致力于定义 $p$ 进 $\GL_2(\Qp)$-Banach 空间 $\Pi_p(f)$ (对任意本征模形式 $f$) 及其最初性质, 然后陈述猜想 \ref{conj:1.1.1}.
第 4 部分包含定理 \ref{thm:1.1.3} 的证明.
最后, 第 5 部分致力于证明定理 \ref{thm:1.1.4}, 然后证明定理 \ref{thm:1.1.2}.
一个附录结束文章, 其中证明了对于定理 \ref{thm:1.1.4} 的证明重要而技术性的一个命题.

\subsection{\mbox{}}

固定 $p$ 为素数.
记 $\overline{\Q}$ 为 $\Q$ 的代数闭包, $\overline{\Qp}$ 为 $\Qp$ 的代数闭包.
固定嵌入 $\overline{\Q} \hookrightarrow \mathbb{C}$ 和 $\overline{\Q} \hookrightarrow \Qp$.
对 $z \in \overline{\Qp}$, $\val(z) \in \Q \cup \{+\infty\}$ 是按 $\val(p) = 1$ 归一化的 $p$ 进赋值, 且 $|z| := p^{-\val(z)} \in \mathbb{R}_{+}$.
记 $\mathbb{A}$ 为 $\Q$ 的 adele, $\mathbb{A}_f := \widehat{\mathbb{Z}} \otimes_{\mathbb{Z}} \Q$ 为有限 adele.
还记 $\widehat{\mathbb{Z}}{}^{p} := \prod_{\ell \neq p} \mathbb{Z}_{\ell}$, 以及 $\mathbb{Z}_{Mp}^{\times} := \varprojlim_n (\mathbb{Z}/Mp^{n}\mathbb{Z})^{\times}$, 其中 $M$ 是满足 $(M,p) = 1$ 的整数.

我们归一化局部类域论的互反律, 把算术 Frobenius 送到一致化子的逆.
记 $\varepsilon$ 为 $\Gal(\Q/\Q)$ 的 $p$ 进分圆特征标, 并注意到, 通过 $p$ 处局部互反律, 若 $z \in \Qp^{\times}$ 则 $\varepsilon(z) = z|z|$.
记 $\nr(x)$ 为 $\Qp^{\times}$ 的将 $p$ 送到 $x$ 的非分歧特征标.

记 $\GL_2$ 为 $\mathbb{Z}$ 上通常的 $2 \times 2$ 可逆矩阵群概形, $B$ (resp.\ $\underline{B}$) 为上 (resp.\ 下) 三角矩阵子群概形, $N$ (resp.\ $\underline{N}$) 为 $B$ (resp.\ $\underline{B}$) 的上 (resp.\ 下) 幺幂矩阵子群概形, $P$ 为 $B$ 的形如 $\begin{pmatrix} * & * \\ 0 & 1 \end{pmatrix}$ 的矩阵子群概形, 而 $T := B \cap \underline{B}$ 为 $\GL_2$ 的对角矩阵子群概形.
对环 $A$ (交换幺环), 记 $\GL_2(A)$ 等为相应的 $A$ 值点群 (我们主要使用 $A = \Qp$ 和 $A = \Zp$).

记 $\mathfrak{gl}_2(\Qp)$, $\mathfrak{b}(\Qp)$, $\mathfrak{t}(\Qp)$, $\mathfrak{n}(\Qp)$ 分别为 $\GL_2(\Qp)$, $B(\Qp)$, $T(\Qp)$, $N(\Qp)$ 的 Lie 代数, 而 $\mathfrak{z}(\Qp)$ 为 $\mathfrak{gl}_2(\Qp)$ 的中心.
设:
\begin{align}
T(\Qp)^{+} &:= \left\{ t \in T(\Qp) \;\middle|\; tN(\Zp)t^{-1} \subset N(\Zp) \right\} = \left\{ \begin{pmatrix} a & 0 \\ 0 & d \end{pmatrix}, \; a,d \in \Qp^{\times}, \; \frac{a}{d} \in \Zp \right\},\\
B(\Qp)^{+} &:= N(\Zp)T(\Qp)^{+} = \left\{ \begin{pmatrix} a & z \\ 0 & d \end{pmatrix}, \; a,d \in \Qp^{\times}, \; \frac{z}{d} \in \Zp \right\}
\end{align}
并记 $X_{-} := \begin{pmatrix} 0 & 0 \\ 1 & 0 \end{pmatrix} \in \mathfrak{gl}_2(\Qp)$.

若 $V$ 是带有 $B(\Qp)$ 线性作用的 $\Qp$-向量空间, $t \in T(\Qp)^{+}$, 且 $N_t := \{n \in N(\Zp) \mid t^{-1}nt \in N(\Zp)\}$, 记 $\pi_t : V^{N(\Zp)} \to V^{N(\Zp)}$ 为如下定义的 Hecke 算子:
\[
\pi_t v := \#(N(\Zp)/N_t)^{-1} \sum_{n \in N(\Zp)/N_t} nt \cdot v.
\]
若 $t \in T(\Zp)$, 则 $\pi_t v = t \cdot v$.
算子 $\pi_t$ 保持 $V^{P(\Zp)}$.

若 $f$ 是权 $k \ge 2$、特征标 $\chi$、对同余子群 $\Gamma_1(N)$ 为新的全纯模形式 (不必尖点), 是 Hecke 算子的特征向量, 且 $L$ 是 $\Qp$ 中包含相应特征值的有限扩张, 记 $\sigma(f) : \Gal(\Q/\Q) \to \GL_2(L)$ 为 Deligne 与 $f$ 相关联的半单 $p$ 进表示.
若 $\ell$ 是素数, 用 $\Frob_{\ell}$ 表示 $\ell$ 处的一个算术 Frobenius.
对 $\ell \nmid Np$, $\sigma(f)(\Frob_{\ell}^{-1})$ 的特征多项式是 $X^2 - a_{\ell}X + \ell^{k-1}\chi(\ell)$, 若 $T_{\ell}f = a_{\ell}f$.
若 $f$ 是尖点的, 表示 $\sigma(f)$ 绝对不可约.
若 $\pi_{p,u}(f)$ 表示由 $f$ 生成的 $\GL_2(\mathbb{A}_f)$ 的不可约光滑可容许表示在 $p$ 处的局部分量 (参见 \cite[\S2.4]{ref:20}), 设:
\[
\pi_p(f) := \pi_{p,u}(f) \otimes |\det|^{\frac{2-k}{2}}.
\]
它是 $\GL_2(\Qp)$ 的不可约光滑表示, 中心特征标为 $|\,|^{2-k}\chi_p^{-1}$, 其中 $\chi_p$ 是把 $\chi$ 看作 $\mathbb{A}_{\Q,f}^{\times}$ 的特征标时其在 $p$ 处的局部分量.
例如, 若 $(p,N) = 1$, 则 $\pi_p(f)$ 是从 $B(\Qp)$ 到 $\GL_2(\Qp)$ 的非分歧特征标 $\nr(p^{-1}\beta_p) \otimes \nr(\alpha_p)$ 的光滑抛物线诱导, 其中 $\alpha_p,\beta_p$ 是 $X^2 - a_pX + p^{k-1}\chi(p)$ 的根.

若 $A$ 是环 (交换幺环), 记 $\Sym^{k-2}A^{2}$ 为 $\GL_2(A)$ 的代数表示, 其底层空间是 $A$-模 $\bigoplus_{0 \le j \le k-2} Az^{j}$, 配备左 $A$-线性作用
\[
(g(P))(z) = (-cz + a)^{k-2}\,P\!\left(\frac{dz-b}{-cz+a}\right),
\]
其中 $g = \begin{pmatrix} a & b \\ c & d \end{pmatrix} \in \GL_2(A)$, $P \in \bigoplus_{0 \le j \le k-2} Az^{j}$.

若 $M$ 是与 $p$ 互素的整数, 记 $\mathbb{T}(M)$ (若 $M$ 固定也简记 $\mathbb{T}$) 为变量 $T_{\ell}$ 和 $S_{\ell}$ ($\ell$ 为不整除 $Mp$ 的素数) 在 $\Zp$ 上生成的多项式代数.
有时把 $\mathbb{T}$-模称为"Hecke 模", 把 (定义在 $L$ 上的) 同态 $\lambda : \mathbb{T} \to L$ 称为"Hecke 特征值系", 其中 $L$ 是 $\Qp$ 中 $\Qp$ 的有限扩张.
若 $\lambda$ 是定义在 $L$ 上的 Hecke 特征值系, $X$ 是 Hecke 模, 记 $X_{L}^{\lambda}$ 为 $X \otimes_{\Zp} L$ 中 $\mathbb{T}$ 经 $\lambda$ 作用的子空间.
若 $\lambda$ 来自 $T_{\ell}$ 和 $S_{\ell}$ 的特征向量模形式 $f$, 也记 $X_{L}^{f} := X_{L}^{\lambda}$.
称定义在 $L$ 上的特征值系 $\lambda$ 是 Eisenstein 的, 若存在连续特征标 $\varepsilon_1,\varepsilon_2 : \mathbb{Z}_{Mp}^{\times} \to L^{\times}$ 使得对一切 $\ell \nmid Mp$ 有 $\lambda(T_{\ell}) = \varepsilon_1(\ell) + \varepsilon_2(\ell)$ 且 $\lambda(\ell S_{\ell}) = \varepsilon_1(\ell)\varepsilon_2(\ell)$.
称 Hecke 模 $X$ 是 Eisenstein 的, 若对 $\Qp$ 的任何有限扩张 $L$, 任何使 $X_{L}^{\lambda} \neq 0$ 的定义在 $L$ 上的特征值系 $\lambda$ 都是 Eisenstein 的.

本文中所有 Banach 空间 $B$ 都是 $p$ 进的, 且满足 $\|B\| \subset |L|$, 其中 $L \subset \Qp$ 是系数域, 它总是 $\Qp$ 的有限扩张.
称带有 $\GL_2(\Qp)$ 左 $L$-线性作用的 Banach 空间 $B$ 为酉 $\GL_2(\Qp)$-Banach 空间, 若对一切 $v \in B$ 映射 $\GL_2(\Qp) \to B$, $g \mapsto gv$ 连续, 且对 $B$ 上某个范数 $\|\,\|$ 的选择, 对一切 $g \in \GL_2(\Qp)$ 和 $v \in B$ 有 $\|gv\| = \|v\|$.
酉 $\GL_2(\Qp)$-Banach 空间称为可容许的 (依 \cite[\S3]{ref:43}), 若其对偶 Banach 空间在 $L \otimes_{\mathcal{O}_L} \mathcal{O}_L[[\GL_2(\Zp)]]$ 上有限型, 其中 $\mathcal{O}_L[[\GL_2(\Zp)]] := \varprojlim_{H} \mathcal{O}_L[\GL_2(\Zp)/H]$, 射影极限取遍 $\GL_2(\Zp)$ 的主同余子群.
